\section{FAIR Data and Open Science}\label{ssc:FAIR-Data}

FAIR (Findable, Accessible, Interoperable, Reusable) is a label for data which are findable in the Web, downloadable and can be used for own purposes. The data in the directory structure described in chapter \ref{ssc:data-for-publication} have to be stored in a FAIR way. This means that data from measurements must be clearly understandable (description of the used device with all adjustments, results with units as ASCII files, etc.) and data from simulations must be reproducible (source code, mesh, boundary conditions, etc.; no commercial programs, but their case file). This method also guarantees the fulfilment of Good Scientific Practice.

For all projects funded by public sponsors (and these are the prevailing case on L-IWT) there hold the principle that all gained data and developed methods have to be provided to the public. Therefore, the directory structure described above can be uploaded to zenodo \url{https://zenodo.org/}, a repository for scientific data, where you get a DOI for your data. But note: PDFs from a publisher can only be stored there if this is permitted explicitly by the journal and must be checked carefully. Usually, it is allowed to provide text and diagrams in an document formated by your own.

Data from other research institutes which are included in a publication, e.g. due to a cooperation, have to be saved FAIR according to our guidelines. This means they provide us their data (stored according to chapter  \ref{ssc:data-for-publication}). If they refuse this, they have to provide at least their data on a public repository.

The reuse of your data by others in the sense of the FAIR principles requires the link to a licence (without a statement the default term is "all rights reserved" which only allows viewing but no further usage).
%see https://academia.stackexchange.com/questions/63139/public-dataset-without-license-what-is-allowed
Therefore, common licences are the 'CC-BY' licence for text and  the '3-clause BSD' license (or 'BSD-3') for software.

The FAIR principles belong to Open Data and is one part within the framework of Open Science, with Open Access of scientific papers and Open Source of software as the three main components. Further information, including the Leibniz Open Science guidelines, can be found on the IWT webpage\footnote{https://www.iwt-bremen.de/en/institute/about-us/open-science}.





\subsection{Publishing Protocols, Data, and Code as Standalone Outputs}
\label{sec:standalone-outputs}

Research projects usually report their outcomes in one or two articles that focus
mainly on the methods and results. The collected data and the scripts written for the
analysis are treated as supplementary material and often stay unpublished. This often
hinders methodological transparency and independent reproducibility, and it limits
further use of the data and code.

IWT therefore encourages treating research methods, datasets, and codebases as separate
research outputs. If possible, each should be published with documentation, ideally as
a short paper, and stored and documented according to these guidelines: the directory
structure in \autoref{ssc:data-for-publication}, metadata in \autoref{ssc:Doku+Metadata}, ELN records with QR codes in \autoref{ssc:ELN}, and licenses as described above. These papers act as guides to the associated
resources. They can be published at different stages of a project, which also shows its
progress. The main benefits are:

\begin{enumerate}
	\item \textbf{Protocol paper or pre-registration.} If you publish your methodology
	and analysis plan before starting the analysis, either as a pre-registration or as a
	protocol paper (also known as registered protocols), you formally state your
	hypotheses or core idea and research plan. You can then get feedback and revise the
	plan if needed. This protects against p-hacking, hypothesizing after the results are
	known (HARKing), and other questionable research practices.
	\item \textbf{Data paper.} A curated dataset with metadata is a valuable resource
	that can often support studies beyond your own. Publishing it with a data paper makes
	it easier for others to reproduce or validate your findings and to reuse the dataset
	in their own work.
	\item \textbf{Software paper.} Publishing your codebase or analysis pipeline with a
	software paper makes your methods and analysis fully transparent and lets others
	adapt your code for their own projects.
	\item \textbf{Citability and visibility.} Publishing these outputs as research papers
	with DOIs, or alongside them, makes them visible and citable milestones. Platforms
	such as Google Scholar, ResearchGate, Scopus, and Semantic Scholar track citations of
	these papers, which they often do not do for data or code repositories alone, even
	when those have DOIs.
\end{enumerate}

Researchers interested in this approach can set up an internal collaboration with the
Digital Technologies (DT) department. DT helps develop a concrete publication plan for
the project and shares the work of preparing and publishing these materials. For
details, contact Prof. Dr. Evgeniya Kabliman.\footnote{%
	\url{https://share.iwt-bremen.de/s/nFxcb6raDN49eYx}
	(Password: \texttt{Open\_2026!})}